\subsection*{Which answer images are available to the reader?}
Four readers of the same pictured passage often proposed the same strip of ink. We trained a selector to assign probability to each reader and added the probabilities of identical strips (Fig.~\ref{fig:span-choice}). On 384 human questions from previously excluded paragraphs, the selector and the generic reader each returned 285 exact images. The selector gained three answers and lost three. It answered 101/192 simplified-script and 184/192 traditional-script questions correctly. The specialist trained on exchanged names returned 178 exact strips and a selector trained with permuted labels returned 225. Performance on the earlier broad sample rose from the specialist's 119 to 181 of 256, below the generic reader's 187. On 32 historical name substitutions, complete constructed statements fell from 17 with generic object selection and 12 with specialized selection to 11 with pooling.

Writing the answers in both requested word orders exposed a different trade-off. The pooled reader and writer passed 33 of the 48 human factual questions, 63 of 128 complete passage exchanges and 34 of 48 reordered-clause tests. The respective retention requirements were 34, 79 and 37. On independent DuIE passages 921/1,200 answers passed both word orders. Thus the gain on broad questions did not preserve all of the specialist's factual writing.

Of 99 fresh direct-answer errors, 15 had an exact reference among the proposals. In the other 84 cases every reader missed it. Widening each reader's list to two, four and eight spans increased the number with an available exact answer from 300 to 319, 337 and 346. The wider search measured availability; selecting reliably among these additional spans remains an empirical question.

